\documentclass[a4paper, 12pt]{article}
\usepackage{amsmath, amsfonts, amsthm, amssymb}
\usepackage{color, fullpage, hyperref}
\usepackage{graphicx}

\setlength{\parindent}{0pt}
\begin{document}
\section{Number System}
\textbf{\underline{Theorem}} Archimedean Property

Let $a,b\in \mathbb{R^+}, \exists n\in\mathbb{N}$, s.t. $na>b$.

\vspace{5mm}

\textbf{\underline{Theorem}} Density of rationals in real numbers
\begin{proof}

    $$\textrm{wts. } \forall c,d\in\mathbb{R},\exists m,n\in\mathbb{N}, ~s.t.~ 
    c<\frac{m}{n}<d$$

    Which is equivalent to $cn < m < dn$.
    By A.P. there exists $n\in\mathbb{N}$ s.t. $\frac{1}{d-c} < n$. So we have
    $$\begin{aligned}
        \frac{1}{d-c} &< n\\
        nc &< nd - 1
    \end{aligned}$$
    Pick an integer $m$, such that $m-1 \leq nc < m$.
    By the inequality on the left hand side, 
    $$\begin{aligned}
        m - 1 &\leq nc\\
        m & \leq nc + 1\\
          &< nd - 1 + 1\\
          &= nd &(*)
    \end{aligned}$$
    Combining with the inequality on the right hand side with $(*)$, the proof
    is completed
    $$\begin{aligned}
        nc < &m < nd\\
        c < &\frac{m}{n} < d
    \end{aligned}$$
\end{proof}
\subsection{Supremum and Infimum}
\textbf{\underline{Definition}} Bouneded

A set $X$ is said to be \textit{bounded} iff. for all $x\in X$, there exists
$l \in \mathbb{R}$, s.t. $x < l$.

\vspace{5mm}
\textbf{\underline{Definition}} Supremum and infimum
Let $X\subseteq \mathbb{R}$ be a set. $x\in\mathbb{R}$ is the \textit{supremum}
of $X$ iff. it satisfies the following:
\begin{enumerate}
    \item $s$ is an upper bound of $X$
    \item for all $\epsilon > 0$, there exists $x\in X$ s.t. $s - \epsilon < x$
\end{enumerate}
Infimum is defined in similar way.

\vspace{5mm}
\textbf{\underline{Theorem }} Existance of square root
$\forall x\in \mathbb{R}, x > 0, \exists y\in \mathbb{R}, y > 0$, s.t. $y^2 = x$

\begin{proof}
    Let $S = \{s\in\mathbb{R^+} :~ s^2 < x\}$.
    Firstly, $S$ is non-empty because $0^2 = 0 < x$. 

    So let $y$ be the supremum of $S$. Want to prove $y^2 = x$.

    Prove by contradiction, and assume $y^2 < x$ or $y^2 > x$.

    \textbf{Case 1: } When $y^2 < x$
    Consider $(y + \frac{1}{n})^2$
    $$\begin{aligned}
        (y + \frac{1}{n})^2 &= y^2 + \frac{2y}{n} + \frac{1}{n^2}\\
                            &< y^2 + \frac{2y}{n} + \frac{1}{n}\\
                            &= y^2 + \frac{2y + 1}{n}   &(*)
    \end{aligned}$$
    By A.P there exists $n_0\in \mathbb{N}$ s.t. 
    $\frac{1}{n} < \frac{x - y^2}{2y + 1}$
    So as continuing $(*)$, we result in 
    $$\begin{aligned}
        (y + \frac{1}{n})^2 &< y^2 + \frac{2y + 1}{n}\\
                            &< y^2 + (2y + 1)\cdot \frac{x - y^2}{2y + 1}\\
                            &= x
    \end{aligned}$$
    This shows $y^2$ is not the supremum of $S$.

    \textbf{Case 2:} When $y^2 > x$
    Consider $(y - \frac{1}{n})^2$, we have
    $$\begin{aligned}
        (y - \frac{1}{n})^2 &= y^2 - \frac{2}{n} + \frac{1}{n^2}\\
                            &> y^2 - \frac{2}{n}  &(**)
    \end{aligned}$$
    Hence, by A.P., there exists $n_0 \in \mathbb{N}$, s.t. 
    $n > \frac{2}{y^2 - x}$.
    So as continuing $(**)$, we result in 
    $$\begin{aligned}
        (y - \frac{1}{n})^2 &> y^2 - \frac{2}{n}\\
                            &> y^2 - 2 \cdot \frac{y^2 - x}{2}\\
                            &= x
    \end{aligned}$$
    This shows $y^2$ is not the supremum of $S$.

    Since $y^2$ neither greater or smaller than $x$, so we have a contradiction.
    Therefore, $y^2 = x$.
\end{proof}

\subsection{Definition of Number System}
\textbf{\underline{Definition }} Natural Numbers

The definition of natural number is also called \textit{Peano Axioms}:
\begin{enumerate}
    \item $1\in \mathbb{N}$
    \item $n\in\mathbb{N} \implies n + 1\in \mathbb{N}$
    \item $\nexists n \in \mathbb{N}$ s.t. $n + 1 = 1$
    \item $\forall n,m\in \mathbb{N}, n + 1 = m + 1 \implies n = m$
    \item Induction property:
        Let $P(n)$ is the property of $n$. If:
        \begin{enumerate}
            \item $P(1)$ is true and
            \item for some $k\in\mathbb{N}, P(k)\implies P(k+1)$
        \end{enumerate}
        Then $P(n)$ is true for all $n\in\mathbb{N}$
\end{enumerate}

\textbf{\underline{Definition}} Equivalence Relation

Let $S$ be set, and for all $a,b\in S$, $\sim$ is an \textit{equivalence relation}
if and only if:
\begin{enumerate}
    \item $a\sim a$
    \item $a\sim b\implies b\sim a$
    \item $a\sim b$ and $b\sim c\sim \implies a\sim c$
\end{enumerate}

\textbf{\underline{Definition }} Equivalence Class

An \textit{equivalence class} is a subset of $A$ that contains all the equivalent
elements.

\vspace{5mm}
\textbf{\underline{Definition }} Integer

Let $A:=\mathbb{N} \times \mathbb{N}$. Define the equivalence relation under
$A$ as:
$$\begin{aligned}
    (n, m) &\sim (n', m') \textrm{ iff.}\\
    n + m' &= n' + m
\end{aligned}$$ 

So $\mathbb{Z}:= A \slash \sim$

\vspace{5mm}
\textbf{\underline{Definition}} Rational numbers

$\mathbb{Q}:= \mathbb{Z}\times (\mathbb{Z} - \{0\})\slash\sim$, where 
$(a, b)\sim (a', b')$ iff. $ab' = a'b$.

\section{Sequence and Series}

\textbf{\underline{Definition }} Cauchy Sequence

A sequence $\{a_n\}$ is said to be \textit{Cauchy} if and only if 
$\forall\epsilon > 0, \exists N > 0 ,$ s.t. 
$m, n > 0 \implies |a_m - a_n| < \epsilon$

\section{Basic Topology on $\mathbb{R}$}
\textbf{\underline{Definition}} Neighbourhood

Let $s\in\mathbb{R}$


\end{document}
